\begin{table}[htbp]
\centering
\caption{Respuestas Agregadas a un Shock de 10\% al Precio del Petróleo}
\label{tab:oil_agg}
\begin{threeparttable}
\begin{tabular}{@{}l c c c@{}}
\toprule
Variable & Impacto & Mín/Máx & Trimestre \\
\midrule
PIB & -0.157 & -0.157 & 1 \\
Inflación IPC & +0.155 & +0.211 & 2 \\
Consumo & -0.187 & -0.201 & 2 \\
Empleo & +0.282 & +0.282 & 1 \\
Tipo de Cambio Real & -0.532 & -0.559 & 2 \\
Balanza Comercial & +0.026 & +0.104 & 3 \\
Tasa de Política & +0.099 & +0.229 & 3 \\
\bottomrule
\end{tabular}
\begin{tablenotes}[flushleft]
\footnotesize
\item \textit{Notas:} PIB, consumo, empleo y TCR en \% de desviación del estado estacionario; inflación y tasa de política en pp anualizados; balanza comercial en pp del PIB. El shock es un aumento único de 10\% en el precio mundial del petróleo $P^{O*}_t$ con persistencia $\rho = 0.9$. Impacto = respuesta en el trimestre 1. Mín/Máx = respuesta extrema en 40 trimestres.
\end{tablenotes}
\end{threeparttable}
\end{table}
